\documentclass[../main.tex]{subfiles}

\begin{document}

\section{Task 1 --- Rigid Launch Vehicle}\label{sec:task1}

\subsection{Rigid model and loop}

Neglecting the bending mode, the plant reduces to the four states
$[\,z,\dot z,\theta,\dot\theta\,]$ and the actuator is taken as ideal:
$W_{TVC}(s)=1$ and no transport delay. The open-loop transfer broken at the TVC
deflection, in the $1+L$ convention, is obtained by closing the static
PD law of Eq.~\eqref{eq:control} around the plant:
\begin{equation}
L(s) = K_{P,\theta}\,G_{\delta\to\theta}(s)
     + K_{D,\theta}\,G_{\delta\to\dot\theta}(s)
     + K_{P,z}\,G_{\delta\to z}(s)
     + K_{D,z}\,G_{\delta\to\dot z}(s),
\end{equation}
where $G_{\delta\to\bullet}$ are the plant channels from $\delta$ to each
measurement. The dominant term is the pitch loop
$L_\theta(s) = (K_{P,\theta}+K_{D,\theta}s)\,K_1/(s^2-A_6)$, whose closed-loop
characteristic polynomial $s^2 + K_1 K_{D,\theta}\,s + (K_1 K_{P,\theta}-A_6)$
is stable only for $K_{P,\theta} > A_6/K_1 = 0.74$ and $K_{D,\theta} > 0$: the
proportional gain must overcome the aerodynamic instability.

\subsection{Tuning}

The pitch gains were tuned (by a small \texttt{fminsearch} wrapper around the
margin computation) to meet the assignment targets, while the lateral gains were
fixed at $K_{P,z}=K_{D,z}=-10^{-3}$. The resulting design is
\begin{equation}
K_{P,\theta} = 1.98,\quad K_{D,\theta} = 1.40,\quad
K_{P,z}=K_{D,z}=-10^{-3},
\end{equation}
with margins $|GM| = 6.0$~dB at $0.34$~rad/s and $|PM| = 30^\circ$ at
$0.21$~rad/s.
In the pole-placement language of the course notes the retained gains place
the dominant pitch pair at $\omega_c = \sqrt{K_1 K_{P,\theta}-A_6} = 2.4$~rad/s
--- inside the typical $1$--$4$~rad/s control band --- with an equivalent
damping $\zeta = K_1 K_{D,\theta}/(2\omega_c) = 1.34$, deliberately above the
$0.7$--$0.8$ pole-placement range because here the damping is driven by the
phase-margin target rather than by a step-response specification.

These relations invert the closed-loop characteristic polynomial: matching
$s^2 + K_1 K_{D,\theta}\,s + (K_1 K_{P,\theta} - A_6)$ to $s^2 + 2\zeta\omega_c s + \omega_c^2$
gives the proportional--derivative design equations
\begin{equation}
K_{P,\theta} = \frac{\omega_c^2 + A_6}{K_1}, \qquad
K_{D,\theta} = \frac{2\zeta\omega_c}{K_1},
\label{eq:poleplace}
\end{equation}
the form tabulated by Barrows and Orr for TVC attitude
control~\cite[Eqs.~(9.2.37)--(9.2.38)]{barrows2020}. The same reference gives a
closed-form estimate of the low-frequency aerodynamic gain margin,
$GM_a \approx 20\log_{10}(K_1 K_{P,\theta}/A_6) = 20\log_{10}(4.565\times1.98/3.382)
\approx 8.5$~dB~\cite[Eq.~(9.2.39)]{barrows2020}; this single-channel figure sits
a couple of dB above the $|GM| = 6.0$~dB returned by \texttt{margin} on the full
four-state loop, the difference being the weak drift channel and the
conditionally-stable low-frequency reading. The gains themselves were obtained
by the \texttt{fminsearch} wrapper minimising the equal-weight margin miss
$(|GM| - 6)^2 + (|PM| - 30)^2$ from the initial guess $[2.0,\,1.4]$, so the
reported $(1.98,\,1.40)$ is the local minimiser of that objective.
The Nichols chart of Figure~\ref{fig:t1nichols} shows the curve threading
between the two critical points $(-180^\circ,0\,\text{dB})$ and
$(-540^\circ,0\,\text{dB})$ --- the geometric signature of conditional stability,
the two crossings being the $360^\circ$ wraps of the single-encirclement Nyquist
condition $N = P_R = 1$ of Section~\ref{sec:intro} --- and clearing the $6$~dB
closed-loop contour; the markers locate the phase and gain crossovers where the
margins are read.

\begin{figure}[H]
\centering
\includegraphics[width=0.62\textwidth]{task1_nichols.png}
\caption{Task 1 --- Nichols chart of the rigid open-loop transfer. The loop is
conditionally stable; the design meets $|GM|\approx 6$~dB and
$|PM|\approx 30^\circ$ (markers at the crossover frequencies).}
\label{fig:t1nichols}
\end{figure}

\subsection{Wind-gust response}

The closed loop is excited by a deterministic ``1-cosine'' wind gust whose
amplitude ($V_g = 6.4$~m/s) is taken from the reference dry-wind dispersion
(\texttt{drywind.mat}, severe regime) at the current altitude, converted to a
wind angle of attack $\alpha_w = v_w/V$ with a peak of $0.39^\circ$. The four
key time histories requested by the assignment are shown in
Figure~\ref{fig:t1gust}: the pitch attitude peaks at $0.139^\circ$, the TVC
deflection at $0.38^\circ$, and the lateral drift builds up smoothly to
$2.5$~m, as expected from the weak drift feedback. The response is well damped
and the actuator effort is modest.

\begin{figure}[H]
\centering
\includegraphics[width=0.85\textwidth]{task1_gust_response.png}
\caption{Task 1 --- response to a severe wind gust: pitch attitude $\theta$,
lateral drift $z$, drift rate $\dot z$ and TVC deflection $\delta$.}
\label{fig:t1gust}
\end{figure}

Since max-$\bar q$ is the load-critical point of the trajectory, the total
angle of attack $\alpha = \theta + \dot z/V + \alpha_w$ and the structural-load
indicator $\bar q\,\alpha$ are also monitored (Figure~\ref{fig:t1alpha}). The
attitude loop pitches the vehicle slightly into the wind
($\theta \to -0.14^\circ$ at the gust peak), so the peak total incidence,
$0.31^\circ$, stays \emph{below} the wind contribution alone --- a mild
load-relieving action. With $\bar q = 81$~kPa the peak load indicator is
$\bar q\,\alpha = 25$~kPa\,deg, the incidence remaining an order of magnitude
below the few-degree load limit typical of a slender launcher, so at this gust
level the drift-minimum--oriented gain choice does not conflict with the
load-minimum requirement.

This drift-versus-load tension is the classical dichotomy of launch-vehicle wind
control~\cite[Sec.~3.1]{greensite_cr826}. The characteristic equation of the
rigid airframe in wind is a cubic with two oscillatory rigid-body roots plus one
real \emph{drift} (flight-path) root; the final-value theorem applied to a step
gust yields a steady drift rate $\dot z_{ss} = -V\,\alpha_w$, independent of the
gains --- which is why the lateral drift of Figure~\ref{fig:t1gust} accumulates
smoothly rather than being nulled. Two limiting philosophies bracket the design:
the \emph{drift-minimum} condition (drift-channel coefficient $B_0 = 0$) holds
the trajectory dispersion down, while the \emph{load-minimum} condition
($K_{P,\theta} = 0$, so the steady incidence $\alpha_{ss} = 0$) minimises
$\bar q\,\alpha$ at the cost of a slow unstable drift mode. The weak negative
drift gains adopted here are a small step toward the drift-minimum end of this
trade.

\begin{figure}[H]
\centering
\includegraphics[width=0.95\textwidth]{task1_alpha_load.png}
\caption{Task 1 --- angle-of-attack budget $\alpha = \theta + \dot z/V +
\alpha_w$ (left) and aerodynamic load indicator $\bar q\,\alpha$ (right) during
the severe gust.}
\label{fig:t1alpha}
\end{figure}

\end{document}
